\section{Conclusiones}
Se estudió la respuesta a distintas señales de un dispositivo desconocido con el fin de modelar su comportamiento con un circuito equivalente. En primer lugar, se sometió el aparato a señales sinusoidales de manera de obtener la impedancia del sistema. En segundo lugar, se utilizaron ondas cuadradas para caracterizar los capacitores en el circuito al ver la carga y descarga de los mismos. 

Para el primer experimento, se obtuvo el diagrama de Nyquist del sistema observando un desplazamiento del sistema en la parte real de la impedancia. A su vez, se observa una semicircunferencia interrumpida por una semirecta y conforme baja la frecuencia la parte imaginaria de la impedancia tiene una asintota vertical. El elemento de fase constante se caracterizó mediante un ajuste de la parte imaginaria en función de la real de la impedancia, de manera que se obtuvo el parametro $n=0.359(6)$ siendo menor al exponente de un elemento de Warburg, indicando que este CPE tiene un comportamiento más resistivo que capacitivo. Finalmente, se propone que el circuito es una conexión de una resistencia en serie $R_s$, con una conexión en paralelo de una resistencia $R_p$ y un capacitor$C_p$ y para explicar el comportamiento de semirecta se introdujo un elemento de fase constante $CPE$ conectado en serie con un capacitor $C_s$. Bajo este modelo, el ajuste obtuvo un parámetro de bondad de $\chi_{\nu}^2=1.965$, indicando que el circuito equivalente propuesto es fiable al dispositivo analizado. 

Para el segundo experimento, observando el diagrama de bode se identifican tres secciones distintivas de la fase del sistema: una zona de crecimiento, luego se mantiene constante, y devuelta aumento. Por ello, se puso al sistema bajo ondas cuadradas en un rango de frecuencias atribuible a cada zona detectada. De esta manera, viendo la atenuación de la señal de salida se ajustó la ecuación de carga y descarga caracteristica de los capacitores. En algunas zonas, no alcanzaba con el comportamiento de un solo capacitor si no que se utilizaron dos. Finalmente, se obtuvieron distintas constantes de tiempo acorde a la región trabajada observando que hay una constante de tiempo que se mantiene entre los [$64$, $85$] \si{\micro\second}, y luego otro que varía entre los [$580$, $3000$] \si{\micro\second}. Estas constantes de tiempo permiten identificar la existencia de distintos capacitores, en el dispositivo analizado se observa que a lo sumo existe un capacitor distinguible, atribuible al primer rango, y luego debido a la variación en el segundo rango pueden existir más. Debido a que el ajuste es bondadoso con 2 capacitores, no se añadieron más al mismo. 

Finalmente, el análisis de espectroscopia de impedancia permitió caracterizar un circuito desconocido, identificando los distintos dispositivos que lo podrían componer como sus valores especificos. Este análisis presentado puede ser extendido a distintos tipos de aparatos tanto inorganicos como órganicos, destacando su importancia en este último.